\chapter{Введение}
% Примечание: я хочу избежать использования ссылок, \term{} и чего-либо подобного здесь. Это просто введение. Все термины определены позже. Давайте будем как можно менее формальными
В области криптографии \textit{доказательства с нулевым разглашением} (zero-knowledge proofs) или \textit{протоколы с нулевым разглашением} (zero-knowledge protocols) — это класс протоколов (protocols), которые позволяют стороне, известной как доказывающий (prover), продемонстрировать истинность утверждения другим сторонам, называемым проверяющими (verifiers), не раскрывая никакой информации, кроме самого факта истинности утверждения. Цель этой книги — дать всестороннее введение в математические основы (mathematical foundations) и реализации (implementations) этих систем доказательств, предназначенное для людей с ограниченным предыдущим опытом в этой области исследований.

Особого значения в этом контексте приобретают \textit{краткие неинтерактивные аргументы знания с нулевым разглашением} (zero-knowledge succinct, non-interactive arguments of knowledge, zk-SNARKs), которые обладают преимуществом: их размер значительно меньше исходных данных, необходимых для установления истинности утверждения, а проверяющему (verifier) достаточно одного сообщения от доказывающего (prover).

С практической точки зрения zk-SNARKs представляют интерес, поскольку они позволяют публично доказать честность вычисления (computation) данных, не раскрывая входные данные (inputs) вычисления, путём передачи краткой транзакции (transaction) проверяющему (verifier), воплощённому в виде смарт-контракта (smart contract) в публичном блокчейне (blockchain). Это облегчает публичную проверку вычислений, повышает масштабируемость (scalability) технологии блокчейн (blockchain) и улучшает конфиденциальность (privacy) транзакций (transactions).

На основе этой взаимосвязи между блокчейнами (blockchains) и zk-SNARKs растущий интерес к технологии блокчейн (blockchain) повысил потребность в более тонком и полном понимании протоколов (protocols) с нулевым разглашением, их реальных приложений и реализаций (implementations), а также разработки стандартов в этой области. Это имеет решающее значение для обеспечения того, чтобы разработчики (developers) могли создавать безопасный и высококачественный код.

Однако тонкости доказательств с нулевым разглашением (zero-knowledge proofs) сложны и требуют глубокого понимания нескольких математических и компьютерно-теоретических дисциплин, а также знакомства с альтернативными моделями вычислений и парадигмами программирования. К сожалению, ресурсы часто разбросаны по блогам, библиотекам на GitHub и математическим статьям, и в результате zk-SNARKs остаются в некоторой степени неуловимыми или ``магическими'', поэтому их иногда называют ``moon math''. Это создаёт барьер для входа и отпугивает разработчиков (developers) от изучения или внедрения их в свои проекты, препятствуя широкому распространению этой технологии и эволюции общества в направлении web3.

``Руководство MoonMath по zk-SNARKs'' призвано изменить это положение; оно разработано специально для людей с ограниченным предыдущим опытом в криптографии (cryptography). Руководство стремится восполнить пробелы в знаниях, предлагая практический подход к объяснению абстрактных понятий с помощью простых вычислений ручкой и бумагой (pen-and-paper). По мере изучения руководства читатели получат понимание математики, лежащей в основе zk-SNARKs, что даст им необходимый фундамент для дальнейшего изучения.

\section{Целевая аудитория}
Основной фокус этой книги — на разработчиков программного обеспечения и смарт-контрактов (smart contract developers), которые стремятся получить глубокое понимание работы zk-SNARKs, чтобы иметь возможность разрабатывать высококачественный и высокозащищённый код, или которые хотят восполнить некоторые пробелы в знаниях. Книга подходит как для начинающих, так и для опытных читателей, поскольку концепции постепенно вводятся структурированным и логичным образом, обеспечивая лёгкость восприятия информации.

Несмотря на то что книга доступна широкому кругу читателей, предполагается, что читатель обладает базовыми знаниями программирования и интересом к логическому мышлению и стратегическому решению задач. Необходим также энтузиазм по отношению к предмету, поскольку детали доказательств с нулевым разглашением (zero-knowledge proofs) могут быть сложными. 


\section{О книге}
Как много математики нужно, чтобы понимать и реализовывать (implement) доказательства с нулевым разглашением (zero-knowledge proofs)? Конечно, ответ зависит от желаемого уровня понимания и требований к безопасности приложения. Возможно реализовать доказательства с нулевым разглашением (zero-knowledge proofs), не понимая лежащей в основе математики; однако чтобы прочитать основополагающую статью, постичь тонкости системы доказательств или разработать безопасные и высококачественные zk-SNARKs, некоторые математические знания необходимы.

Без прочной математической подготовки тот, кто заинтересован в изучении концепций доказательств с нулевым разглашением (zero-knowledge proofs), но никогда не видел и не имел дела, скажем, с простым полем или эллиптической кривой, может быстро почувствовать себя перегруженным. Дело не столько в сложности необходимой математики, сколько в огромном количестве технического жаргона, незнакомых терминов и непонятных символов, которые быстро делают текст нечитаемым, даже если сами концепции на самом деле не так уж сложны. В результате читатель может либо потерять интерес, либо подхватить некоторые несвязные кусочки знаний, которые в худшем случае приводят к незрелым и небезопасным реализациям (implementations). 

Значительная часть этой книги посвящена всестороннему объяснению математических основ (mathematical foundations), необходимых для понимания фундаментальных концепций, лежащих в основе разработки zk-SNARKs. Читателям, не знакомым с основами теории чисел и эллиптическими кривыми, мы настоятельно рекомендуем тщательно изучить соответствующие главы (chapters), пока они не смогут решить как минимум несколько упражнений (exercises) в каждой главе (chapter). Настоятельно рекомендуется целенаправленно внимательно разбирать примеры (examples).

Книга начинается с очень базового уровня и предполагает лишь наличие предварительных знаний фундаментальных концепций, таких как арифметика целых чисел в объёме средней школы. Затем она постепенно показывает, что существуют числа и математические структуры, которые, хотя поначалу кажутся совсем непохожими на то, что изучалось в школе, на более глубоком уровне на самом деле аналогичны. Это иллюстрируется разнообразными примерами (examples) по всей книге.

Важно подчеркнуть, что математика, представленная в этой книге, неформальна, неполна и оптимизирована для облегчения эффективного понимания концепций с нулевым разглашением. Выбор в пользу включения только минимально необходимой теории с приоритетом обилия числовых примеров (examples) сделан намеренно. Мы считаем, что такой неформальный, основанный на примерах (example-driven) подход облегчает усвоение материала для начинающих.

Начинающему рекомендуется сначала вычислить простой учебный zk-SNARK с помощью ручки и бумаги (pen-and-paper), прежде чем приступать к разработке высокозащищённых zk-SNARKs для реального мира. Однако для вывода таких учебных примеров (examples) необходима некоторая математическая подготовка. Поэтому эта книга ставит целью направить неопытного читателя к основным математическим концепциям, предоставляя упражнения (exercises), которые предназначены для самостоятельной работы. Каждый раздел (section) включает список постепенно усложняющихся упражнений (exercises), чтобы помочь запомнить и применить концепции.


\section{Как читать эту книгу}
Руководство MoonMath предназначено для всестороннего введения в темы, актуальные для начинающих в математике и криптографии (cryptography). Поэтому существует несколько способов чтения книги. Наиболее прямолинейный подход — следовать главам (chapters) в линейном порядке. Этот метод рекомендуется читателям с ограниченными предварительными знаниями математики и криптографии (cryptography). Книга начинается с фундаментальных концепций, таких как натуральные числа, простые числа и операции над этими множествами в различных арифметиках. Затем книга переходит к алгебраическим структурам, включая группы, кольца, простые поля и эллиптические кривые.

На протяжении всей книги вводятся примеры (examples), которые постепенно расширяются за счёт включения новых знаний из последующих глав (chapters). Такой поэтапный подход позволяет развивать простые, но полноценные криптографические системы, которые можно вычислять вручную, чтобы дать подробную иллюстрацию каждого шага.

Читателям, интересующимся пониманием эллиптических кривых в контексте систем доказательств с нулевым разглашением, отправной точкой может служить введение в кривую $BLS6_6$ в \secname{} \ref{BLS6}. Кривая $BLS6_6$ была специально разработана для вычислений вручную, поскольку размер криптографических эллиптических кривых часто делает невозможными такие вычисления. Это кривая, дружественная к спариванию (pairing-friendly), обладающая всеми необходимыми свойствами для выполнения вычислений на основе спаривания без помощи компьютера, что может помочь прояснить тонкости системы. Кроме того, в книге приводится вывод кривой \curvename{Tiny-jubjub} \ref{TJJ13-twisted-edwards}, которая может использоваться для вычислений EdDSA в схемах над $BLS6_6$.

Читателям, заинтересованным в построении простого zk-SNARK с нуля с помощью ручки и бумаги (pen-and-paper), стоит начать с примеров (examples), связанных с задачей 3-факторизации. В \examplename{} \ref{ex:3-factorization} мы вводим задачу 3-факторизации как утверждение в формальном языке. Если это слишком абстрактно, читатель может начать с \examplename{} \ref{ex:3-fac-zk-circuit}, где мы описываем задачу 3-факторизации как алгебраическую схему. В \examplename{} \ref{ex:3-fac-zk-circuit_2} мы выполняем схему. В \examplename{} \ref{ex:L-3fac-zk} мы вводим в задачу концепции instance (публичный вход) и witness (свидетель), чтобы впоследствии достичь различных уровней нулевого разглашения. В \examplename{} \ref{ex:3-factorization-r1cs} мы преобразуем схему в соответствующую систему ограничений ранга 1 (Rank-1 Constraint System), а в \examplename{} \ref{ex:3-fac-R1CS-constr-proof} вычисляем конструктивное доказательство для задачи. В \examplename{} \ref{ex:3-fac-QAP} мы преобразуем эту систему ограничений в квадратичную арифметическую программу (Quadratic Arithmetic Program) и показываем, как конструктивные доказательства превращаются в задачи о делимости многочленов. В \examplename{} \ref{ex:3-fac-groth-16-params} мы используем результат этих примеров (examples) для вывода Groth16 zk-SNARK для задачи 3-факторизации. В \ref{def:groth16-crs} мы вычисляем ключ доказывающего (prover) и ключ проверяющего (verifier). В \examplename{} \ref{3-fac-snark-compute} мы вычисляем zk-SNARK, а в \examplename{} \ref{3-fac-snark-verifier} — проверяем этот zk-SNARK. В \examplename{} \ref{3-fac-snark-simulator} мы показываем, как моделировать доказательства. 

Читателям, желающим реализовать (implement) zk-SNARK на практическом языке программирования, следует начать с нашей реализации (implementation) на Circom задачи 3-факторизации, как описано в \ref{ex:3-fac-circom}. Вывод соответствующего набора параметров Groth\_16 изложен в \ref{ex:3-fac-groth-16-params-circom}, а вычисление связанной общей эталонной строки (Common Reference String) рассмотрено в \ref{ex:3-fac-groth-16-setup-circom}. В \ref{ex:3-fac-groth-16-prover-circom} мы демонстрируем генерацию доказательства для случайно выбранных входных значений, которое затем проверяется в \ref{ex:3-fac-groth-16-verifier-circom}. Эта иллюстрация может служить отправной точкой для более глубокого понимания лежащих в основе механизмов.

Лица, стремящиеся углубить своё понимание процесса, посредством которого высокоуровневые программы компилируются в представления, пригодные для анализа системами доказательств с нулевым разглашением, могут извлечь пользу из изучения главы (Chapter) \ref{chap:circuit-compilers}, в которой мы развиваем игрушечный язык, оснащённый ``мозг-компилятором'' (brain-compiler), преобразующим высокоуровневый код в графические представления схем. Ключевые представления, такие как R1CS, подробно рассмотрены в разделе (Section) \ref{sec:R1CS}, а фундаментальные принципы конструктивных доказательств, свидетелей (witnesses) и входных данных (instances) изложены в \ref{sec:formal-languages}.



\section{Вклад}Учитывая значимость и стремительное развитие области систем доказательств с нулевым разглашением (zero-knowledge proofing systems), всестороннее освещение соответствующих тем в Руководстве MoonMath представляет собой серьёзную задачу. Поэтому вклад сообщества крайне ценится.

Если вы хотите внести вклад в развитие Руководства MoonMath, мы призываем вас направлять ваши адаптированные материалы или оригинальный контент в Least Authority. Тем, кто заинтересован в существенном вкладе в Руководство MoonMath, мы предлагаем связаться непосредственно с Least Authority по адресу mmm@leastauthority.com. Для получения дополнительной информации, пожалуйста, обратитесь к нашей лицензии.

%остальное закомментировано

\begin{comment}
\section{Зоопарк доказательств с нулевым разглашением}

{Сначала список систем доказательств с нулевым разглашением:

\begin{enumerate}
	\item Pinocchio (2013): \href{https://eprint.iacr.org/2013/279.pdf}{{Paper}}
	\begin{itemize}[label={--}]
		\item Примечания: trusted setup
	\end{itemize}

	\item BCGTV (2013): \href{https://eprint.iacr.org/2013/507.pdf}{{Paper}}
	\begin{itemize}[label={--}]
		\item Примечания: trusted setup, implementation
	\end{itemize}

	\item BCTV (2013): \href{https://eprint.iacr.org/2013/879.pdf}{{Paper}}
	\begin{itemize}[label={--}]
		\item Примечания: trusted setup, implementation
	\end{itemize}

	\item Groth16 (2016): 	\href{https://eprint.iacr.org/2016/260.pdf}{Paper}
	\begin{itemize}[label={--}]
		\item Примечания: trusted setup
		\item Другие ресурсы: \href{https://www.gakonst.com/zksummit2019.pdf}{Talk in 2019 by Georgios Konstantopoulos}
	\end{itemize}

	\item GM17 (207): 	\href{https://eprint.iacr.org/2017/540.pdf}{Paper}
	\begin{itemize}[label={--}]
		\item Примечания: trusted setup
		\item Другие ресурсы: позже \href{https://eprint.iacr.org/2018/187}{Simulation extractability in ROM, 2018}
	\end{itemize}

	\item Bulletproofs (2017): \href{https://eprint.iacr.org/2017/1066.pdf}{Paper}
	\begin{itemize}[label={--}]
		\item Примечания: no trusted setup
		\item Другие ресурсы: \href{https://eprint.iacr.org/2016/263.pdf}{Polynomial Commitment Scheme on DL, 2016} и \href{https://www.iacr.org/archive/asiacrypt2010/6477178/6477178.pdf}{KZG10, Polynomial Commitment Scheme on Pairings, 2010}
	\end{itemize}

	\item Ligero (2017): \href{https://acmccs.github.io/papers/p2087-amesA.pdf}{Paper}
	\begin{itemize}[label={--}]
		\item Примечания: no trusted setup
		\item Другие ресурсы: 
	\end{itemize}

	\item Hyrax (2017): \href{https://eprint.iacr.org/2017/1132.pdf}{Paper}
	\begin{itemize}[label={--}]
		\item Примечания: no trusted setup
		\item Другие ресурсы: 
	\end{itemize}

	\item STARKs (2018): \href{https://eprint.iacr.org/2018/046.pdf}{Paper}
	\begin{itemize}[label={--}]
		\item Примечания: no trusted setup 
		\item Другие ресурсы: 
	\end{itemize}

	\item Aurora (2018): \href{https://eprint.iacr.org/2018/828.pdf}{Paper}
	\begin{itemize}[label={--}]
		\item Примечания: transparent SNARK
		\item Другие ресурсы:
	\end{itemize}

	\item Sonic (2019): \href{https://eprint.iacr.org/2019/099.pdf}{Paper}
	\begin{itemize}[label={--}]
		\item Примечания: SNORK — SNARK с универсальной и обновляемой доверенной настройкой (universal and updateable trusted setup), основанный на PCS
		\item Другие ресурсы: \href{https://www.benthamsgaze.org/2019/02/07/introducing-sonic-a-practical-zk-snark-with-a-nearly-trustless-setup/}{Blog post by Mary Maller from 2019} и \href{https://eprint.iacr.org/2018/280}{work on updateable and universal setup from 2018}
	\end{itemize}

	\item Libra (2019): \href{https://eprint.iacr.org/2019/317}{Paper}
	\begin{itemize}[label={--}]
		\item Примечания: trusted setup
		\item Другие ресурсы:
	\end{itemize}

	\item Spartan (2019): \href{https://eprint.iacr.org/2019/550.pdf}{Paper}
	\begin{itemize}[label={--}]
		\item Примечания: transparent SNARK
		\item Другие ресурсы:
	\end{itemize}

	\item PLONK (2019): \href{https://eprint.iacr.org/2019/953.pdf}{Paper}
	\begin{itemize}[label={--}]
		\item Примечания: SNORK, основанный на PCS
		\item Другие ресурсы: \href{https://www.plonk.cafe/t/welcome-to-discussion-of-plonk-related-research/24}{Discussion on Plonk systems} и \href{https://github.com/Fluidex/awesome-plonk}{Awesome Plonk list}
	\end{itemize}

	\item Halo (2019): \href{https://eprint.iacr.org/2019/1021}{Paper}
	\begin{itemize}[label={--}]
		\item Примечания: no trusted setup, основанный на PCS, рекурсивный
		\item Другие ресурсы: 
	\end{itemize}

	\item Marlin (2019): \href{https://eprint.iacr.org/2019/1047.pdf}{Paper}
	\begin{itemize}[label={--}]
		\item Примечания: SNORK, основанный на PCS
		\item Другие ресурсы: \href{https://github.com/arkworks-rs/marlin}{Rust Github}
	\end{itemize}

	\item Fractal (2019): \href{https://eprint.iacr.org/2019/1076.pdf}{Paper}
	\begin{itemize}[label={--}]
		\item Примечания: рекурсивный, transparent SNARK
		\item Другие ресурсы: 
	\end{itemize}

	\item SuperSonic (2019): \href{https://eprint.iacr.org/2019/1229.pdf}{Paper}
	\begin{itemize}[label={--}]
		\item Примечания: transparent SNARK, основанный на PCS
		\item Другие ресурсы: \href{https://eprint.iacr.org/2021/358}{Attack on DARK compiler in 2021}
	\end{itemize}

	\item Redshift (2019): \href{https://eprint.iacr.org/2019/1400.pdf}{Paper}
	\begin{itemize}[label={--}]
		\item Примечания: SNORK, основанный на PCS
		\item Другие ресурсы: 
	\end{itemize}



\end{enumerate}

\textbf{Другие ресурсы по зоопарку: } \href{https://github.com/matter-labs/awesome-zero-knowledge-proofs}{Awesome ZKP list on Github}, \href{https://zkp.science/}{ZKP community} with the \href{https://docs.zkproof.org/reference.pdf}{reference document}

}

\paragraph{Список дел}
\begin{itemize}
	\item Составить таблицу для времени доказывающего (prover), времени проверяющего (verifier) и размера доказательства
	\item Подумать о категориях — \textit{Достигнутые цели}: trusted setup или нет, постквантовый или нет, \dots
	\item Подумать о категориях — \textit{Математическая подготовка}: polynomial commitment scheme, \dots
	\item \dots пока мы обсуждаем пункты выше, следует также обсудить общую нотацию/язык для всего этого. (Например, transparent SNARK/no trusted setup/STARK)
\end{itemize}

\paragraph{Вопросы, которые следует осветить при написании}
\begin{itemize}
	\item Сделать исторический обзор ``открытия'' различных систем ZKP
	\item Помочь читателю понять, какая статья основана на каком результате и т.д. — дерево публикаций!
	\item Помочь читателю понять различную терминологию, например SNARK/SNORK/STARK, PCS, R1CS, updateable, universal, $\dots$
	\item Помочь читателю понять математические предпосылки — и что это означает для зоопарка.
	\item Куда пойдёт развитие/эволюция? Какие узкие места?
\end{itemize}

\vspace*{1em}
{\footnotesize
\hspace*{-1em}\textbf{Другие темы, на которые я наткнулся при составлении этого списка}
\begin{itemize}
	\item Vector commitments: \url{https://eprint.iacr.org/2020/527.pdf}
	\item Snarkl: \url{http://ace.cs.ohio.edu/~gstewart/papers/snaarkl.pdf}
	\item Virgo?: \url{https://people.eecs.berkeley.edu/~kubitron/courses/cs262a-F19/projects/reports/project5_report_ver2.pdf}
\end{itemize} 
}

\end{comment}
